\section{Dimensions de la grille de déplacement du robot en coordonnées polaires}
Les dimensions horizontales $X$ et verticales $Y$ d'un carreau de la grille de dé\-pla\-ce\-ment du
\robot\ peuvent être définies en coordonnées cartésiennes par les commandes
{\idt lgX(}$X${\idt )} et {\idt lgY(}$Y${\idt )} comme indiqué à la figure \ref{fig:lgrt} ci-dessous.

\begin{figure}[h]
 \includevspacePostScript[4.0cm]{LgrtExplique}{0.95}{0.95}
 \parbox{13.7cm}{\caption{Grille de déplacement du \protect\robot\ définie par}\label{fig:lgrt}} 
 \centering{{\idt lgX(}$X${\idt );\ lgY(}$Y${\idt )}\ ou\ {\idt lgRT(}$R${\idt ,} $T${\idt )}\Large\strut} 
\end{figure}

On peut également définir les dimensions d'un carreau du quadrillage en {\em coordonnées
polaires} \index{Coordonnées polaires}par la commande
{\idt lgRT(}$R${\idt ,} $T${\idt )} en donnant la longueur $R$ de la diagonale de ce carreau mesurée
en pixels et
l'angle $T$  exprimé en degrés que fait cette diagonale avec une horizontale
comme illustré à la figure \ref{fig:lgrt} ci-dessus. La commande {\idt lgRT} a deux
paramètres, qui sont des nombres entiers ou rationnels, comme dans {\idt lgRT(50, 70)} ou
{\idt lgRT(50.5, 22.5)}.

\begin{exemple}[Triangle isocèle]{\rm
Le programme ci-dessous dessine un triangle isocèle ayant deux angles égaux
à 70 degrés et deux côtés égaux à 50 pixels~:

{\raggedright 
 %\vspace{\baselineskip}
 \vbox to 0pt{\includevspacePostScript[4cm]{TriangleIsocele}{0.95}{0.95}}%
 \input python/TriangleIsocele.tex
 }
\noindent $\Box$ }
\end{exemple}

\section{Rotations}
\index{Rotation}La commande {\idt rt(}{\em t}{\idt )} fait effectuer au \robot\ une \underline{r}o\underline{t}ation de {\em t} degrés,
{\em t} étant une expression entière ou rationnelle (par exemple {\idt rt(15)}, {\idt rt(-90)} ou {\idt rt(22.5)}).
Si la valeur de l'expression {\em t} est positive alors le pivotement sur place se fait dans le sens 
{\em contraire} des aiguilles d'une montre\index{Sens trigonométrique}\footnote{\ Sens contraire des aiguilles d'une montre que les mathématiciens
appellent {\em sens trigonométrique}.}. Si la valeur de l'expression {\em t} est négative alors le pivotement sur 
place se fait dans le sens des ai\-guil\-les d'une montre. Cette commande ne change pas la position du \robot\ mais elle change
son orientation. La longueur du segment tracé par le prochain {\idt av()} reste également inchangée. Pour ce faire la taille du 
carreau du quadrillage doit être convenablement modifiée. Ceci est illustré sur la figure \ref{fig:rtexp} 
(page \pageref{fig:rtexp}) qui montre  un carreau du quadrillage en pointillés avant et après des rotations successives de
15 degrés (le \robot\ partant à chaque fois du coin inférieur gauche du carreau et la rotation étant suivie d'un {\idt av()} pour montrer 
que la longueur du segment tracé ne change pas). Sur cette figure, les 
schémas successifs sont disposés de haut en bas et de gauche à droite.
\begin{figure}[htbp]
 \includevspacePostScript[17.7cm]{RTexplique}{0.95}{0.95}
 \caption{Rotations et déplacements successifs {\idt rt(15);\ av()}}\label{fig:rtexp}
\end{figure}

Quand on utilise la commande {\idt rt}, nous conseillons d'éviter les commandes {\idt phtd()}, 
{\idt pqtd()},  {\idt p3htd()}, \ldots, {\idt pne()}, {\idt pse()}, \ldots en les remplaçant par {\idt rt} avec un paramètre
approprié (seules les commandes {\idt pn()}, {\idt ps()}, \ldots sont utiles pour les positionnements absolus). 
En utilisant exclusivement la commande {\idt  rt}, il n'est pas nécessaire de tenir 
compte des changements de repères  qu'elle induit. 
\begin{exemple}[Pentagone]{\rm
Le programme ci-dessous dessine un pentagone en répétant cinq fois le tracé d'un segment suivi
d'une rotation de $(360 / 5) = 72$ degrés~:

 {\raggedright 
 \vbox to 0pt{\includevspacePostScript[6cm]{Pentagone_2}{0.95}{0.95}}%
 \input python/Pentagone_2.tex
 \vspace{1.6cm}% le listing Python est plus court que la figure placee a cote
 }
\noindent $\Box$ }\end{exemple}
\begin{exemple}[Rosace de pentagones]{\rm
La rosace ci-dessous est obtenue par superposition de dix pentagones tournés successivement de 36 degrés~:

 {\raggedright 
% \vspace{\baselineskip}
 \vbox to 0pt{\includevspacePostScript[6cm]{RosaceDePentagones}{0.95}{0.95}}%
 \input python/RosaceDePentagones.tex
 }
\noindent $\Box$ }\end{exemple}

\section{Translations}
\index{Translation}
Quelle que soit l'orientation du \robot, l'expression  {\idt valLgTl()} de type {\idt float} a pour valeur la longueur
 (mesurée en pixels) d'un segment qui serait tracé par un {\idt av()}. 
 
 Pour tracer un segment d'une longueur différente, on peut utiliser la commande {\idt tl(}$E${\idt )} qui 
 prend en paramètre une expression $E$ (entière ou rationnelle) dont la valeur est la 
 longueur de ce segment 
 comptée en pixels. Si cette valeur est positive, le segment est tracé en avançant dans la 
 direction courante du \robot. Si cette valeur est négative, le tracé se fait en reculant. Le \robot\
 est déplacé à l'extrémité du segment et ne change pas son orientation. Les dimensions
 d'un carreau du quadrillage sont changées pour permettre ce déplacement (si le crayon est levé)
 ou ce tracé (si le crayon est baissé). Par conséquent, le segment peut être prolongé
 par un segment d'égale longueur en utilisant une commande {\idt av()}.
 
 Par exemple, la séquence de commandes {\idt tl(20);\ rt(90);\ av();\ 
rt(90);\ av();\ rt(90);\ av()} trace un carré de côté
20 pixels.
\begin{exemple}[Grecque polygonale]{\rm
La frise polygonale ci-dessous est ob\-te\-nue en remplaçant le tracé
d'un côté du polygone par celui de la grecque~:

\vspace{\baselineskip}

\includevspacePostScript[6.5cm]{GrecquePolygonale}{0.95}{0.95}
\input python/GrecquePolygonale.tex
\noindent $\Box$ }\end{exemple}
\noindent \rule{137mm}{0.1mm} %-----------------------------------------------------------
%%% EXOS lgrt
\begin{exercices}{Écrire des programmes \python\  utilisant la commande {\idt lgRT}
pour réaliser les dessins suivants~:

\tiret un triangle équilatéral$\:$;

\tiret le pentagone, les croix ou la toile d'araignée de la figure \ref{fig:PtCrCrTa} (page \pageref{fig:PtCrCrTa})$\:$;

\tiret les motifs géométriques d'après Horemis \cite{horemis} de la figure \ref{fig:horemis} (page \pageref{fig:horemis})$\:$;

\tiret les pavages coloriés de la figure \ref{fig:pavtrilo} (page \pageref{fig:pavtrilo}).

Les très nombreux pavages donnés par Grünbaum et Shephard \cite{grunbaum} sont la source
d'innombrables exercices.
\noindent $\Box$}\end{exercices}
\begin{figure}[htbp]
 \includevspacePostScript[3.1cm]{PtCrCrTa}{0.95}{0.95}
 \caption{Pentagone, croix et toile d'araignée}\label{fig:PtCrCrTa}
\end{figure}
\begin{figure}[htbp]
 \includevspacePostScript[6.0cm]{Horemis}{0.95}{0.95}
 \caption{Motifs géométriques d'après Horemis}\label{fig:horemis}
\end{figure}
\begin{figure}[htbp]
 \includevspacePostScript[3.2cm]{EtoilesORCVSpi}{0.95}{0.95}
 \caption{Étoiles, octogone décoré et spirolatérale}\label{fig:etocsp}
\end{figure}
\begin{figure}[htbp]
 \includevspacePostScript[10.2cm]{PavagesTriLos}{0.95}{0.95}
 \caption{Pavages de triangles, losanges et carrés}\label{fig:pavtrilo}
\end{figure}
\begin{figure}[htbp]
 \includevspacePostScript[4.4cm]{RosacesPolygonales}{0.95}{0.95}
 \caption{Rosaces polygonales}\label{fig:rospol}
\end{figure}
%%% EXOS rt
\begin{exercices}{Écrire des programmes \python\  utilisant la commande {\idt rt}
pour réaliser les dessins suivants~:

\tiret les étoiles, l'octogone décoré ou la courbe spirolatérale d'ordre 3 à 135 degrés de la figure 
         \ref{fig:etocsp} (page \pageref{fig:etocsp}), cette dernière étant obtenue en avançant une, puis deux, puis
		 trois fois, en tournant à chaque fois de 135 degrés et en répétant ces déplacements jusqu'à ce que
		 la courbe soit fermée$\:$;

\tiret les polygones réguliers de la figure 
         \ref{fig:polyreg} (page \pageref{fig:polyreg}), le nombre de côtés étant un paramètre de la procédure
		 de dessin$\:$;

\tiret les étoiles de la figure 
         \ref{fig:etoiles} (page \pageref{fig:etoiles}), le nombre de branches étant un paramètre de la procédure
		 de dessin$\:$;

\tiret le rapporteur de la figure 
         \ref{fig:rapport} (page \pageref{fig:rapport}), en utilisant la commande
		 {\idt ecrire(entierEnChaine(}{\em e}{\idt ))} qui écrit la valeur de l'expression entière {\em e}
		 à droite de la position courante du \robot\ (l'irrégularité des graduations
		 est due aux erreurs d'arrondi en coordonnées d'écran entières qui ne permettent pas de
		 représenter des droites de pentes quelconques)$\:$;

\tiret le pavage du plan avec des hexagones ou des étoiles de la figure 
         \ref{fig:pavhexeto} (page \pageref{fig:pavhexeto})$\:$;

\tiret le pavage à la Escher \cite{escher1} du disque et du triangle de la figure 
         \ref{fig:pavdistri} (page \pageref{fig:pavdistri})$\:$;

\tiret le pavage du disque utilisant le pavé de Voderberg \cite{voderberg} de la figure 
         \ref{fig:voder1} (page \pageref{fig:voder1})$\:$;

\tiret les spirales utilisant le pavé de Voderberg \cite{voderberg} de la 
         figure \ref{fig:voder2} (page \pageref{fig:voder2}).

\noindent $\Box$}\end{exercices}
\begin{figure}[htbp]
 \includevspacePostScript[7.3cm]{PolygonesReguliers}{0.95}{0.95}
 \caption{Polygones réguliers}\label{fig:polyreg}
\end{figure}
\begin{figure}[htbp]
 \includevspacePostScript[1.9cm]{Etoiles}{0.95}{0.95}
 \caption{Étoiles}\label{fig:etoiles}
\end{figure}
\begin{figure}[htbp]
 \includevspacePostScript[6.8cm]{Rapporteur}{0.95}{0.95}
 \caption{Rapporteur gradué en degrés}\label{fig:rapport}
\end{figure}
\begin{figure}[htbp]
 \includevspacePostScript[6.0cm]{PavageHexagonesEtoiles}{0.95}{0.95}
 \caption{Pavage du plan avec des hexagones et des étoiles}\label{fig:pavhexeto}
\end{figure}
\begin{figure}[htbp]
 \includevspacePostScript[9.7cm]{PavageDisqueTriangle}{0.95}{0.95}
 \caption{Pavage à la Escher du disque et du triangle}\label{fig:pavdistri}
\end{figure}
\begin{figure}[htbp]
 \includevspacePostScript[6.5cm]{Voderberg1}{0.95}{0.95}
 \caption{Pavé et pavage du disque de Voderberg}\label{fig:voder1}
\end{figure}
\begin{figure}[htbp]
 \includevspacePostScript[6.2cm]{Voderberg2}{0.95}{0.95}
 \caption{Spirales de Voderberg}\label{fig:voder2}
\end{figure}
%%% EXOS developpements de solides
\begin{exercices}{Écrire des programmes \python\  utilisant les commandes {\idt lgRT} et {\idt rt} 
pour réaliser les développements des solides suivants~:

\tiret solides de Platon (cube, tétraèdre, octaèdre, dodécaèdre et iso\-ca\-è\-dre
ré\-gu\-liers) de la figure \ref{fig:platon} (page \pageref{fig:platon})$\:$;

\tiret les grand et petit rhombicuoctaèdres et le tétraèdre tronqué de la figure 
\ref{fig:devsol} (page \pageref{fig:devsol}).

Les développements des solides donnés par Cundy et Rollett \cite{cundy} 
sont la source de nombreux exercices.
\noindent $\Box$}\end{exercices}
\begin{figure}[htbp]
 \includevspacePostScript[5.5cm]{SolidesDePlaton}{0.95}{0.95}
 \caption{Développements des solides de Platon}
 \label{fig:platon}
\end{figure}
\begin{figure}[htbp]
\includevspacePostScript[6.0cm]{DevelopSolides}{0.95}{0.95}
\caption{Développements des grand et petit rhombicuoctaèdres}\label{fig:devsol}
\centering{et du tétraèdre tronqué}
\end{figure}

%%% EXOS rt + tl + tracer
\begin{exercices}{Écrire des programmes \python\  utilisant les commandes {\idt rt}
et {\idt tl} pour réaliser les dessins suivants~:

\tiret les rosaces polygonales de la figure \ref{fig:rospol} (page \pageref{fig:rospol})$\:$;

\tiret les ``jolygones'' de la figure \ref{fig:jgones} (page \pageref{fig:jgones}) obtenus en traçant des
segments successifs dont la longueur diminue à chaque fois de 5 \%, deux segments successifs faisant
un angle de $\theta$ degrés$\:$; $\theta$ étant choisi respectivement égal à 15, 60, 72, 76, 89, 92,
118, 175 et 144 degrés
pour les jolygones 
%de la figure \ref{fig:jgones} (page \pageref{fig:jgones}) 
considérés de gauche
à droite et de haut en bas$\:$;

\tiret les courbes planes de la figure \ref{fig:courbpo} (page \pageref{fig:courbpo})
définies par une équation polaire, comme suit~:
 \begin{list}{$\scriptscriptstyle\bullet$}{
}
   \item le limaçon de Pascal, $\rho = a + b \cos{\theta}$ avec
   $a \neq 0$, $b \neq 0$ et $\theta \in [0, 2\pi]$ (sur l'exemple, $a = 40$ et $b = 80$),
   \item le rhodone\ae\ à trois pétales, $\rho = a \cos{n \theta}$ avec 
    $a = 60$, $n = 3$ et $\theta \in [0, \pi]$,
   \item la courbe de Moritz, $\rho = a \cos{\frac{9}{10} \theta}$ avec 
    $a = 60$ et $\theta \in [-10 \pi, 10 \pi]$$\:$;
   \end{list}

\tiret les courbes planes définies par les équations cartésiennes $y = \sin{x}, x \in [-2\pi, 2\pi]$
  (figure \ref{fig:sin} ci-dessous) et  $y = \cos{x}, x \in [-2\pi, 2\pi]$ en utilisant les commandes {\idt deplacerReel} 
  et {\idt tracerReel}. La commande {\idt tracerReel(x, y)},
  où {\idt x} et {\idt y} sont des expressions rationnelles,
  trace un segment (si le crayon est baissé) de la position courante du
  \robot\ au point de coordonnées cartésiennes $(x,\ y)$ et déplace le \robot\ dans
  cette nouvelle position  $(x,\ y)$. La commande {\idt deplacerReel(x, y)} fait de
  même sans tracer. Les fonctions {\idt sin} et {\idt cos} et la constante {\idt pi}
  sont expliquées dans l'addenda du chapitre \ref{ch:rec}.
   De nombreux autres exemples de tracés de courbes sont donnés dans \cite{LT}.

\noindent $\Box$}\end{exercices}
\begin{figure}[htbp]
 \includevspacePostScript[4.7cm]{Sinus}{0.95}{0.95}
 \caption{Courbe cartésienne $y = \sin{x}, x \in [-2\pi, 2\pi]$}\label{fig:sin}
\end{figure}
\begin{figure}[htbp]
 \includevspacePostScript[12.0cm]{Jolygones}{0.95}{0.95}
 \caption{Jolygones à 15, 60, 72, 76, 89, 92, 118, 175 et 144 degrés}\label{fig:jgones}
\end{figure}
\begin{figure}[htbp]
 \includevspacePostScript[5.3cm]{CourbesPolaires}{0.95}{0.95}
 \caption{Limaçon de Pascal, rhodone\ae\ à trois pétales, courbe de Moritz}\label{fig:courbpo}
\end{figure}
%%% CORRIGES lgrt
\begin{corrige}{
Nous donnons le programme de dessin de la toile d'araignée de la figure \ref{fig:PtCrCrTa} 
(page \pageref{fig:PtCrCrTa}), les autres programmes se trouvant dans le dossier
{\idt livre/\allowbreak 20\_grille\_polaire}~:
\input python/ToileDaraignee.tex
\noindent $\Box$}\end{corrige}
%%% CORRIGES rt
\begin{corrige}{
Nous donnons la procédure de dessin d'un polygone régulier de $n$ côtés (figure \ref{fig:polyreg}, 
page \pageref{fig:polyreg}), les autres programmes se trouvant  dans le dossier
{\idt livre/\allowbreak 20\_grille\_polaire}~:
\input python/PolygoneRegulier.tex
\noindent $\Box$}\end{corrige}
%%% CORRIGES developpements de solides
\begin{corrige}{
Nous donnons le programme de dessin du développement du té\-traè\-dre tron\-qué, les autres 
programmes se trouvant  dans le dossier {\idt livre/\allowbreak 20\_grille\_polaire}~:
\input python/TetraedreTronque.tex
\noindent $\Box$}\end{corrige}
%%% CORRIGES rt + tl
\begin{corrige}{
Nous donnons le programme de dessin du jolygone à 175 degrés, les autres jolygones sont obtenus en changeant la
valeur de la constante {\idt angle}. Les autres programmes se trouvent  dans le dossier
{\idt livre/\allowbreak 20\_grille\_polaire}.
\input python/Jolygone175.tex
\noindent $\Box$}\end{corrige}

\clearpage

%%%%%%%%%%%%%%%%%%%%%%%%%%%%%%%%%%%%%%%%%%%%%%%%%%%%%%%%%%%%%%%%%%%%%%%%%%%
\noindent \framebox[\textwidth]{
\begin{minipage}{\myparboxwidth}
\setlength{\parindent}{2em}
\small \sf
\vspace{\abovedisplayskip}
\section*{Addenda}
La commande {\idt lgRT(r, t)} définit la longueur {\idt r} de la diagonale d'un carreau du
quadrillage (exprimée en pixels) et l'angle {\idt t} exprimé en degrés que fait cette diagonale 
avec une horizontale.  Elle est équivalente à~:
\begin{center}
{\idt lgX(r * cos(t / 180 * pi))}\\
{\idt lgY(r * sin(t / 180 * pi))} 
\end{center}
mais évite l'emploi des fonctions trigonométriques.
L'angle {\idt t}, compris entre 0 et 90,  définit l'orientation nord-est qui est confondue avec l'est quand {\idt t} vaut {\idt 0}
et avec le nord quand {\idt t} vaut {\idt 90}.

La commande {\idt rt(t)} fait pivoter le \robot\ sur place d'un angle de {\idt t} degrés à gauche,
dans le sens trigonométrique ou sens inverse des aiguilles d'une montre, si {\idt t} $> 0$, et à droite, dans le sens des aiguilles
d'une montre, si {\idt t} $ < 0$. La commande {\idt tl(l)} fait avancer (si {\idt l} $> 0$)
ou reculer (si {\idt l} $< 0$) le \robot\ tout droit de {\idt l} fois la taille d'un côté d'un pixel
de l'écran.

Ces commandes {\idt rt} et {\idt tl} sont directement inspirées du graphisme ``tortue'' de 
{\sc logo}. Elles changent le repère et l'orientation du \robot\ pour rester compatibles
avec l'emploi de coordonnées cartésiennes ou polaires.

Le langage {\sc logo} fut conçu par Seymour Papert à la fin des années 1970 \cite{logo}.
Il présente de nombreuses similitudes avec {\sc lisp}
\cite{lisp}, en particulier l'interprétation de données représentant des instructions,
la portée dynamique des identificateurs et la structure de liste.
L'interprétation des données comme des programmes implique quelques subtilités dans
l'écriture des programmes (comme l'emploi de la fonction {\idt quote}, \cite{wilson}) qui ne sont pas
à la portée de tous les débutants en informatique.
La portée dynamique des identificateurs \cite{wilson}
est largement abandonnée et la structure de liste
est presque toujours complétée par des tableaux dans les  {\sc lisp} modernes comme
{\sc common lisp} qui suivent en cela les langages descendant d'{\sc algol 60} 
\cite{algol60} comme \pascal\ ou {\sc ada}. Le langage \python\ offre
lui aussi à la fois les listes et les tableaux, n'interprète pas ses données comme
des instructions (sauf si on le lui demande explicitement) et utilise la portée statique des identificateurs.
 
\vspace{\abovedisplayskip}
\end{minipage}
}
%%%%%%%%%%%%%%%%%%%%%%%%%%%%%%%%%%%%%%%%%%%%%%%%%%%%%%%%%%%%%%%%%%%%%%%%%%%% 
\newpage                 % Si on termine le chapitre par une page blanche
\thispagestyle{empty}
